% Generated draft: requires mathematical and source review.
\documentclass[11pt]{article}
\usepackage[margin=25mm]{geometry}
\usepackage{fontspec}
\setmainfont{Latin Modern Roman}
\usepackage{amsmath,amssymb,mathtools,unicode-math}
\setmathfont{Latin Modern Math}
\usepackage{enumitem,needspace,xurl,hyperref}
\setlength{\parindent}{0pt}
\setlength{\parskip}{5pt}
\setlength{\emergencystretch}{3em}
\newcommand{\N}{\mathbb{N}}
\newcommand{\Z}{\mathbb{Z}}
\newcommand{\Q}{\mathbb{Q}}
\newcommand{\R}{\mathbb{R}}
\newcommand{\C}{\mathbb{C}}
\newcommand{\F}{\mathbb{F}}
\newcommand{\A}{\mathbb{A}}
\newcommand{\PP}{\mathbb{P}}
\newcommand{\E}{\mathbb{E}}
\newcommand{\eps}{\varepsilon}
\newcommand{\dd}{\,\mathrm{d}}
\begin{document}
\section*{UnsolvedMath ID: 11100042}
\textbf{Source problem code:} catalogue entry
\subsection{Definitions and mathematical statement}
Let $\mathcal C=\operatorname{Ho}(\operatorname{Sp})$ be the stable homotopy category. An endomorphism of its identity functor is a family $\eta_X:X\to X$ natural in all spectra: for every $f:X\to Y$, $\eta_Yf=f\eta_X$. Addition and composition make these transformations a ring $Z(\mathcal C)$. The excerpt does not specify whether transformations must also commute with suspension or be enriched/homotopy-coherent; those are additional conventions.

Evaluation on the sphere gives a split ring map
\[
 e:Z(\mathcal C)\to[\mathbb S,\mathbb S]\cong\Z,
 \qquad n\longmapsto(n\operatorname{id}_X)_X
\]
as its section. Axiomatic stable homotopy item 10 asks to determine this ring, especially whether $\ker e=0$, so every transformation is an integer multiple of the identity.

A self-map $u:X\to X$ is phantom if $u\circ a=0$ for every map $a:F\to X$ from a finite spectrum. Freyd's generating hypothesis says a map between finite spectra inducing zero on all stable homotopy groups is zero. The source says that, assuming this hypothesis, a kernel transformation is a natural family of self-phantoms, and asks why such a family must vanish and whether the hypothesis is needed. To justify the reduction one must in particular specify the suspension compatibility ensuring evaluation on all sphere suspensions is zero, rather than only on $\mathbb S$.

The registry summary claims the ring is $\Z$ without the generating hypothesis, attributing this to Oka and Mike Hopkins and an unspecified nilpotence/rigidity argument. No exact theorem or reference is supplied. A computation of homotopy-coherent transformations in an enhanced category would not alone settle the ordinary-category formulation; the notion of transformation must match. The registry's open metadata and solved summary are therefore retained as an unresolved source/status discrepancy.
\subsection{Short English statement}
Determine all natural self-transformations of the identity on the stable homotopy category, and whether sphere evaluation has any phantom kernel. Match the ordinary, suspension-compatible or enhanced notion before accepting the supplied claim that the ring is $\Z$.
\subsection{Sources}
\begin{itemize}
\item Local registry: \url{data/unsolvedmath/problems.json}, permanent UnsolvedMath ID $11100042$, exact source code \texttt{catalogue entry}, attributed to Mark Hovey.
\item \emph{Hovey's Algebraic Topology Problem List}, Axiomatic stable homotopy item 10: \url{https://www-users.cse.umn.edu/~tlawson/hovey/}. Supplied access date: 2026-07-29; extraction: html-ordered-list; 65 top-level items in nine category pages.
\item The registry mentions Oka, Mike Hopkins, an unspecified independent write-up, and a cryptic ``Cauchy'' usage as support for a solved claim. No title, identifier, precise category or theorem statement is supplied for these attributions.
\end{itemize}
\end{document}
